\documentclass[11pt,a4paper]{letter}
\usepackage[margin=1in]{geometry}
\usepackage[utf8]{inputenc}
\usepackage[T1]{fontenc}

\signature{Tushar Pandey \\ Texas A\&M University \\ \texttt{tusharp@tamu.edu}}
\address{Tushar Pandey \\ Texas A\&M University \\ College Station, TX, USA}

\begin{document}

\begin{letter}{
  The Editor \\
  \textit{Quantum Science and Technology} \\
  IOP Publishing
}

\opening{Dear Editor,}

I am pleased to submit the manuscript ``Capacity, Not Barren Plateaus:
Diagnosing Variational Quantum Training Failure on Chaotic Dynamics, and
a Fixed-Reservoir Resolution'' for consideration as a research paper in
\textit{Quantum Science and Technology}. The work is original, has not
been published elsewhere, and is not under consideration at any other
journal.

The paper makes a diagnostic contribution that I believe is essential
reading for the quantum-machine-learning sub-field. Variational
quantum-circuit failures on hard learning tasks at NISQ scale are
routinely attributed to barren plateaus --- exponentially vanishing
gradients in deep parameterised circuits --- and a substantial volume
of recent work proposes mitigations targeted at this specific failure
mode. We demonstrate, on a benchmark chaotic-dynamics task at the
relevant hardware scale, that this attribution is empirically and
theoretically wrong. The McClean barren-plateau bound for the relevant
ansatz class predicts gradient magnitudes that are four orders of
magnitude below those we measure. The dominant failure mode is
capacity-limited expressivity rather than gradient vanishing, and
mitigations targeted at the gradient-vanishing regime should not be
expected to help.

This diagnostic finding has lasting consequences. First, it constrains
which mitigation strategies are appropriate at NISQ scale, redirecting
effort toward architectural rather than gradient-side interventions.
Second, the diagnostic methodology --- the direct comparison of
measured gradient norms against the McClean threshold for the ansatz
class --- generalises to any variational quantum-machine-learning
study and should be standard practice. Third, the finding clarifies
what an architectural alternative needs to bypass: not gradients per
se, but the underlying capacity limit. We illustrate this with a
fully specified Quantum Reservoir Computing pipeline that achieves
$81\%$ lower mean-squared error and trains approximately $52{,}000
\times$ faster on matched quantum resources, with the same
architecture generalising across three canonical chaotic systems
without per-system retuning.

The paper additionally contributes an empirical scaling sweep on
Lorenz, $q \in \{5, \dots, 10\}$ against the standard ESN baseline at
$N_{\mathrm{ESN}} = 500$, showing that the naive feature-dimension
parity crossover does not occur: QRC test MSE degrades monotonically
with qubit count at fixed prediction-target dimension because features
dilute relative to a fixed-dimensional signal. The correct scaling
axis is matched-complexity, in which qubit count and target dimension
grow together; this is the subject of a separate companion manuscript
in preparation. The QST paper also contributes a windowed-variational
quantum-circuit (WVQC) ablation that closes the information-asymmetry
concern in the QRC--QPINN comparison. The WVQC
ablation, in which a variational circuit is given the same windowed
input as QRC and trained by Adam under the same protocol as QPINN,
recovers only approximately $32\%$ of the QRC--QPINN gap on matched
seeds, isolating the fixed-reservoir architectural choice as the
dominant factor and extending the capacity-limit diagnosis to the
information-symmetric setting.

All experimental code, data, and figures are released publicly at the
repository linked in the manuscript, supporting reproducibility and
extension by other groups.

I have intentionally restructured the paper so that the diagnostic
finding leads, with the architectural comparison serving as
supporting evidence rather than as the primary claim. I believe this
ordering reflects the lasting contribution and is appropriate for
the journal's selectivity criterion.

I would be grateful for your consideration. I am happy to suggest
reviewers if helpful.

\closing{Sincerely,}

\end{letter}

\end{document}
